\documentclass[11pt, a4paper]{article}
\usepackage[utf8]{inputenc}
\usepackage[margin=1in]{geometry}
\usepackage{listings}
\usepackage{xcolor}
\usepackage{tabularx}

% Configure SQL code block styling
\lstset{
    basicstyle=\ttfamily\small,
    keywordstyle=\color{blue}\bfseries,
    stringstyle=\color{red!70!black},
    commentstyle=\color{green!60!black}\itshape,
    frame=single,
    backgroundcolor=\color{gray!5},
    breaklines=true,
    showstringspaces=false
}

\title{Assignment4 Lab Report}
\author{Student Name}
\date{\today}

\begin{document}

\maketitle

\section{Problem8}

\subsection*{Schema Creation}
\textbf{Schema Requirements:}
\begin{itemize}
    \item Customer (\underline{cust\_id}, cust\_name, annual\_revenue, cust\_type) \\
    \textit{Cust\_id must be between 100 and 10,000. Cust\_type must be `MANUFACTURER', `WHOLESALER', or `RETAILER'.}
    \item Shipment (\underline{shipment\_no}, cust\_id, weight, truck\_no, destination, ship\_date) \\
    \textit{Foreign keys: cust\_id references customer on delete cascade, truck\_no references truck on delete set null, destination references city. Weight must be under 1000.}
    \item Truck (\underline{truck\_no}, driver\_name)
    \item City (\underline{city\_name}, population)
\end{itemize}

\begin{lstlisting}[language=SQL]
-- DDL Schema Creation
CREATE TABLE IF NOT EXISTS Customer (
  cust_id INTEGER PRIMARY KEY CHECK (
    cust_id >= 100
    AND cust_id <= 10000
  ),
  cust_name VARCHAR(30),
  annual_revenue NUMERIC(10, 2),
  cust_type VARCHAR(30) CHECK (
    cust_type IN ('MANUFACTURER', 'WHOLESALER', 'RETAILER')
  )
);

CREATE TABLE IF NOT EXISTS Truck (
  truck_no INTEGER PRIMARY KEY,
  driver_name VARCHAR(30)
);

CREATE TABLE IF NOT EXISTS City (
  city_name VARCHAR(30) PRIMARY KEY,
  population INTEGER
);

CREATE TABLE IF NOT EXISTS Shipment (
  shipment_no INTEGER PRIMARY KEY,
  cust_id INTEGER REFERENCES Customer (cust_id) ON DELETE CASCADE,
  weight INTEGER CHECK (weight < 1000),
  truck_no INTEGER REFERENCES Truck (truck_no) ON DELETE SET NULL,
  destination VARCHAR(30) REFERENCES City (city_name)
);
\end{lstlisting}
\newpage

\subsection*{Q1: Give names of customer who have sent packages (shipments) to Kolkata, Chennai and Mumbai.}
\begin{lstlisting}[language=SQL]
SELECT
  cust_name
FROM
  customer
  JOIN shipment ON customer.cust_id = shipment.cust_id
WHERE
  destination in ('Kolkata', 'Chennai', 'Mumbai');
\end{lstlisting}
\vspace{0.3cm}
\textbf{Output:}\vspace{0.2cm}
\begin{center}
\begin{tabular}{|l|}
\hline
\textbf{cust\_name} \\ \hline
Acme Corp \\ \hline
Acme Corp \\ \hline
Acme Corp \\ \hline
Echo Supplies \\ \hline
Echo Supplies \\ \hline
Echo Supplies \\ \hline
Halo Tech \\ \hline
\end{tabular}
\end{center}
\newpage
\subsection*{Q2: List the names of the driver who have delivered shipments weighing over 200 pounds.}
\begin{lstlisting}[language=SQL]
SELECT
  driver_name,
  weight
FROM
  truck
  JOIN shipment on truck.truck_no = shipment.truck_no
WHERE
  shipment.weight > 200;
\end{lstlisting}
\vspace{0.3cm}
\textbf{Output:}\vspace{0.2cm}
\begin{center}
\begin{tabular}{|l|l|}
\hline
\textbf{driver\_name} & \textbf{weight} \\ \hline
IQBAL & 300 \\ \hline
MANISH & 210 \\ \hline
DEV & 500 \\ \hline
RAJ & 250 \\ \hline
AMIT & 300 \\ \hline
IQBAL & 999 \\ \hline
IQBAL & 750 \\ \hline
\end{tabular}
\end{center}
\newpage
\subsection*{Q3: Retrieve the maximum and minimum weights of the shipments. Rename the output as Max\_Weight and Min\_Weight respectively.}
\begin{lstlisting}[language=SQL]
SELECT
  MAX(weight) as Max_Weight,
  MIN(weight) as Min_Weight
FROM
  shipment;
\end{lstlisting}
\vspace{0.3cm}
\textbf{Output:}\vspace{0.2cm}
\begin{center}
\begin{tabular}{|l|l|}
\hline
\textbf{max\_weight} & \textbf{min\_weight} \\ \hline
999 & 10 \\ \hline
\end{tabular}
\end{center}
\newpage
\subsection*{Q4: For each customer, what is the average weight of package sent by the customer?}
\begin{lstlisting}[language=SQL]
SELECT
  cust_name,
  ROUND(AVG(weight), 2)
FROM
  customer
  JOIN shipment on customer.cust_id = shipment.cust_id
GROUP BY
  customer.cust_id
ORDER BY
  AVG(weight) DESC;
\end{lstlisting}
\vspace{0.3cm}
\textbf{Output:}\vspace{0.2cm}
\begin{center}
\begin{tabular}{|l|l|}
\hline
\textbf{cust\_name} & \textbf{round} \\ \hline
Jupiter Mills & 625.00 \\ \hline
Gamma Traders & 544.50 \\ \hline
Indus Retail & 210.00 \\ \hline
Beta Goods & 200.00 \\ \hline
Echo Supplies & 180.00 \\ \hline
Acme Corp & 150.00 \\ \hline
Delta Dynamics & 120.00 \\ \hline
Charlie Retail & 116.67 \\ \hline
Foxtrot Mart & 80.00 \\ \hline
Halo Tech & 52.50 \\ \hline
\end{tabular}
\end{center}
\newpage
\subsection*{Q5: List the names and populations of cities that have received a shipment weighing over 100 pounds.}
\begin{lstlisting}[language=SQL]
SELECT
  city_name,
  population,
  weight
FROM
  city
  JOIN shipment ON shipment.destination = city.city_name
WHERE
  weight > 100
ORDER BY
  weight ASC;
\end{lstlisting}
\vspace{0.3cm}
\textbf{Output:}\vspace{0.2cm}
\begin{center}
\begin{tabular}{|l|l|l|}
\hline
\textbf{city\_name} & \textbf{population} & \textbf{weight} \\ \hline
Mumbai & 12000000 & 150 \\ \hline
Kolkata & 4500000 & 150 \\ \hline
Chennai & 7000000 & 150 \\ \hline
Pune & 800000 & 150 \\ \hline
Delhi & 11000000 & 150 \\ \hline
Bangalore & 8000000 & 150 \\ \hline
Hyderabad & 6900000 & 150 \\ \hline
Jaipur & 3000000 & 150 \\ \hline
Surat & 4400000 & 150 \\ \hline
Kochi & 600000 & 150 \\ \hline
Pune & 800000 & 210 \\ \hline
Chennai & 7000000 & 250 \\ \hline
Pune & 800000 & 300 \\ \hline
Kolkata & 4500000 & 300 \\ \hline
Pune & 800000 & 500 \\ \hline
Kochi & 600000 & 750 \\ \hline
Delhi & 11000000 & 999 \\ \hline
\end{tabular}
\end{center}
\newpage
\subsection*{Q6: List cities that have received shipments from every customer.}
\begin{lstlisting}[language=SQL]
SELECT
  city.city_name
FROM
  city
  JOIN shipment on shipment.destination = city.city_name
GROUP BY
  city.city_name
HAVING
  COUNT(DISTINCT shipment.cust_id) = (
    SELECT
      COUNT(cust_id)
    FROM
      customer
  );
\end{lstlisting}
\vspace{0.3cm}
\textbf{Output:}\vspace{0.2cm}
\begin{center}
\begin{tabular}{|l|}
\hline
\textbf{city\_name} \\ \hline
Pune \\ \hline
\end{tabular}
\end{center}
\newpage
\subsection*{Q7: For each city, what is the maximum weight of a package sent to that city?}
\begin{lstlisting}[language=SQL]
SELECT
  city.city_name,
  MAX(shipment.weight)
FROM
  city
  JOIN shipment on city.city_name = shipment.destination
GROUP BY
  city.city_name;
\end{lstlisting}
\vspace{0.3cm}
\textbf{Output:}\vspace{0.2cm}
\begin{center}
\begin{tabular}{|l|l|}
\hline
\textbf{city\_name} & \textbf{max} \\ \hline
Kochi & 750 \\ \hline
Hyderabad & 150 \\ \hline
Kolkata & 300 \\ \hline
Pune & 500 \\ \hline
Jaipur & 150 \\ \hline
Mumbai & 150 \\ \hline
Delhi & 999 \\ \hline
Surat & 150 \\ \hline
Chennai & 250 \\ \hline
Bangalore & 150 \\ \hline
\end{tabular}
\end{center}
\newpage
\subsection*{Q8: List the name and annual revenue of customers whose shipments have been delivered by truck driver ‘IQBAL’.}
\begin{lstlisting}[language=SQL]
SELECT DISTINCT
  customer.cust_name,
  customer.annual_revenue
FROM
  customer
  JOIN shipment on customer.cust_id = shipment.cust_id
  JOIN truck on shipment.truck_no = truck.truck_no
WHERE
  driver_name = 'IQBAL';
\end{lstlisting}
\vspace{0.3cm}
\textbf{Output:}\vspace{0.2cm}
\begin{center}
\begin{tabular}{|l|l|}
\hline
\textbf{cust\_name} & \textbf{annual\_revenue} \\ \hline
Beta Goods & 300000.00 \\ \hline
Echo Supplies & 200000.00 \\ \hline
Gamma Traders & 450000.00 \\ \hline
Jupiter Mills & 600000.00 \\ \hline
\end{tabular}
\end{center}
\newpage
\subsection*{Q9: List drivers who have delivered shipments to every city.}
\begin{lstlisting}[language=SQL]
SELECT
  truck.driver_name
FROM
  truck
  JOIN shipment on shipment.truck_no = truck.truck_no
GROUP BY
  truck.truck_no
HAVING
  COUNT(DISTINCT shipment.destination) = (
    SELECT
      COUNT(city_name)
    FROM
      city
  );
\end{lstlisting}
\vspace{0.3cm}
\textbf{Output:}\vspace{0.2cm}
\begin{center}
\begin{tabular}{|l|}
\hline
\textbf{driver\_name} \\ \hline
SAM \\ \hline
\end{tabular}
\end{center}
\newpage
\subsection*{Q10: For each city, with population over 1 million, what is the minimum weight of a package sent to that city.}
\begin{lstlisting}[language=SQL]
SELECT
  city.city_name,
  MIN(shipment.weight)
FROM
  city
  JOIN shipment on shipment.destination = city.city_name
WHERE
  city.population > 1000000
GROUP BY
  city.city_name
\end{lstlisting}
\vspace{0.3cm}
\textbf{Output:}\vspace{0.2cm}
\begin{center}
\begin{tabular}{|l|l|}
\hline
\textbf{city\_name} & \textbf{min} \\ \hline
Mumbai & 10 \\ \hline
Jaipur & 150 \\ \hline
Delhi & 150 \\ \hline
Surat & 150 \\ \hline
Chennai & 150 \\ \hline
Hyderabad & 150 \\ \hline
Bangalore & 150 \\ \hline
Kolkata & 150 \\ \hline
\end{tabular}
\end{center}
\newpage

\section{Problem9}

\subsection*{Schema Creation}
\textbf{Schema Requirements:}
\begin{itemize}
    \item EMP (\underline{EMPNO}, ENAME, JOB, MGR, HIREDATE, SAL, COMM, DEPTNO) \\
    \textit{EMPNO must be between 7000 and 8000. ENAME must not exceed 10 characters. JOB must be in (`Clerk', `Salesman', `Manager', `Analyst', `President'). MGR is the manager's EMPNO. COMM must be under 1500 and defaults to 0.}
    \item DEPT (\underline{DEPTNO}, DNAME, LOC) \\
    \textit{DEPTNO must start with `D'. DNAME must be `Accounting' or `Sales' or `Research' or `Operations'.}
\end{itemize}

\begin{lstlisting}[language=SQL]
-- DDL Schema Creation
DROP TABLE IF EXISTS EMP CASCADE;

DROP TABLE IF EXISTS DEPT CASCADE;

CREATE TABLE IF NOT EXISTS dept (
  deptno VARCHAR(10) PRIMARY KEY CHECK (DEPTNO LIKE 'D%'),
  dname VARCHAR(20) CHECK (
    dname IN ('Accounting', 'Sales', 'Research', 'Operations')
  ),
  loc VARCHAR(30)
);

CREATE TABLE IF NOT EXISTS emp (
  empno INTEGER PRIMARY KEY CHECK (
    empno >= 7000
    AND empno <= 8000
  ),
  ename VARCHAR(10),
  job VARCHAR(15) CHECK (
    job in (
      'Clerk',
      'Salesman',
      'Manager',
      'Analyst',
      'President'
    )
  ),
  mgr INTEGER REFERENCES EMP (empno) ON DELETE SET NULL,
  hiredate DATE,
  sal NUMERIC(10, 2),
  comm NUMERIC(10, 2) DEFAULT 0 CHECK (comm < 1500),
  deptno VARCHAR(10) REFERENCES DEPT (DEPTNO)
);
\end{lstlisting}
\newpage

\subsection*{Q1: Display the difference between highest and lowest salary of each department in descending order. Label the column as “Difference”.}
\begin{lstlisting}[language=SQL]
SELECT
  dname,
  MAX(sal) - MIN(sal) AS Difference
FROM
  dept
  JOIN emp on emp.deptno = dept.deptno
GROUP BY
  dname
ORDER BY
  MAX(sal) - MIN(sal) DESC;
\end{lstlisting}
\vspace{0.3cm}
\textbf{Output:}\vspace{0.2cm}
\begin{center}
\begin{tabular}{|l|l|}
\hline
\textbf{dname} & \textbf{difference} \\ \hline
Accounting & 3700.00 \\ \hline
Research & 2200.00 \\ \hline
Sales & 1900.00 \\ \hline
\end{tabular}
\end{center}
\newpage
\subsection*{Q2: List all the employees’ employee number and name along with their immediate managers’ employee number and name.}
\begin{lstlisting}[language=SQL]
SELECT
  e1.empno,
  e1.ename,
  e2.empno as manager_id,
  e2.ename as manager_name
FROM
  emp e1
  LEFT JOIN emp e2 on e1.mgr = e2.empno;
\end{lstlisting}
\vspace{0.3cm}
\textbf{Output:}\vspace{0.2cm}
\begin{center}
\begin{tabular}{|l|l|l|l|}
\hline
\textbf{empno} & \textbf{ename} & \textbf{manager\_id} & \textbf{manager\_name} \\ \hline
7839 & Andrew & NULL & NULL \\ \hline
7566 & JONES & 7839 & Andrew \\ \hline
7698 & BLAKE & 7839 & Andrew \\ \hline
7782 & CLARK & 7839 & Andrew \\ \hline
7788 & SCOTT & 7566 & JONES \\ \hline
7902 & FORD & 7566 & JONES \\ \hline
7499 & ALLEN & 7698 & BLAKE \\ \hline
7521 & WARD & 7698 & BLAKE \\ \hline
7654 & MARTIN & 7698 & BLAKE \\ \hline
7369 & SMITH & 7902 & FORD \\ \hline
7876 & ADAMS & 7788 & SCOTT \\ \hline
7900 & JAMES & 7698 & BLAKE \\ \hline
7934 & MILLER & 7782 & CLARK \\ \hline
\end{tabular}
\end{center}
\newpage
\subsection*{Q3: Create a query that will display the total number of employees and the total number of employees who were hired only in 2020. Give the column headings as “TOTAL” and “TOTAL\_2020” respectively.}
\begin{lstlisting}[language=SQL]
SELECT
  (
    SELECT
      COUNT(*)
    FROM
      emp
  ) AS TOTAL,
  (
    SELECT
      COUNT(*)
    FROM
      emp
    WHERE
      EXTRACT(
        YEAR
        FROM
          hiredate
      ) = 2020
  ) as TOTAL_2020;

\end{lstlisting}
\vspace{0.3cm}
\textbf{Output:}\vspace{0.2cm}
\begin{center}
\begin{tabular}{|l|l|}
\hline
\textbf{total} & \textbf{total\_2020} \\ \hline
13 & 8 \\ \hline
\end{tabular}
\end{center}
\newpage
\subsection*{Q4: Display the manager number and the salary of the lowest paid employee under that manager. Exclude anyone whose manager is not known. Exclude any group where the minimum salaryis less than 1000. Sort the output in descending order of salary.}
\begin{lstlisting}[language=SQL]
SELECT
  mgr as manager_id,
  MIN(sal) as min_sal
FROM
  emp
WHERE
  mgr IS NOT NULL
GROUP BY
  mgr
HAVING
  MIN(sal) >= 1000
ORDER BY
  min_sal DESC;
\end{lstlisting}
\vspace{0.3cm}
\textbf{Output:}\vspace{0.2cm}
\begin{center}
\begin{tabular}{|l|l|}
\hline
\textbf{manager\_id} & \textbf{min\_sal} \\ \hline
7566 & 3000.00 \\ \hline
7839 & 2450.00 \\ \hline
7782 & 1300.00 \\ \hline
7788 & 1100.00 \\ \hline
\end{tabular}
\end{center}
\newpage
\subsection*{Q5: Assume that there are some departments where no employee is assigned. Now, write
a query to display the department name, location name, number of employees, and the
average salary for all the employees in that department. Label the columns as
“DNAME”, “LOCATION”, “NUMBER OF PEOPLE”, and “AVERAGE SALARY”
respectively. Round the averge salary to two decimal places. The outcome of the
query must include the details of the departments where no employee is assigned and
in that case the “AVERAGE SALARY” for that department is to be displayed as 0
(zero).}
\begin{lstlisting}[language=SQL]
SELECT
  dname,
  loc AS LOCATION,
  COUNT(empno) as "NUMBER OF PEOPLE",
  ROUND(COALESCE(AVG(sal), 0), 2) as "AVERAGE SALARY"
FROM
  dept d
  LEFT JOIN emp e ON d.deptno = e.deptno
GROUP BY
  dname,
  loc
ORDER BY
  "AVERAGE SALARY" DESC;
\end{lstlisting}
\vspace{0.3cm}
\textbf{Output:}\vspace{0.2cm}
\begin{center}
\begin{tabular}{|l|l|l|l|}
\hline
\textbf{dname} & \textbf{location} & \textbf{NUMBER OF PEOPLE} & \textbf{AVERAGE SALARY} \\ \hline
Accounting & NEW YORK & 3 & 2916.67 \\ \hline
Research & DALLAS & 5 & 2175.00 \\ \hline
Sales & CHICAGO & 5 & 1580.00 \\ \hline
Operations & BOSTON & 0 & 0.00 \\ \hline
\end{tabular}
\end{center}
\newpage

\end{document}